\section{Verifier Construction: Moving the Target into the Fast Loop}
\label{sec:verifier-construction}

So far we have largely treated the feedback structure as given.  An optimizer receives
a proxy every \(\tauP\), while more independent target evidence arrives after
\(\tauT\).  But this assumption hides one of the most important ways that difficult
problems become tractable: the feedback architecture itself can be changed.

A useful intervention is therefore not only
\[
    \text{generate better candidates},
\]
but
\[
    \boxed{\text{make mistakes observable sooner and more faithfully}.}
\]

We call this \emph{verifier construction}.  The term includes formal verifiers, tests,
instruments, experiments, surrogate measurements, monitoring systems, and
institutions that move target-relevant information into a shorter feedback loop.

\subsection{Verification is constructed, not merely discovered}

Many domains now treated as ``objectively verifiable'' became so only after substantial
technical and institutional work.  A physical quantity requires an operational
measurement procedure.  An experimental result requires controlled comparisons,
calibration, and a way to distinguish signal from noise.  A software property requires
a specification and a test or proof obligation.  A clinical claim requires an endpoint
and a study design.

Thus the apparent binary distinction
\[
    \text{verifiable} \quad\text{vs.}\quad \text{unverifiable}
\]
is often better represented as
\[
    \text{weakly observable}
    \rightarrow
    \text{instrumented}
    \rightarrow
    \text{repeatably measurable}
    \rightarrow
    \text{cheaply testable}.
\]

Scientific instruments are especially clear examples.  An instrument can turn a
latent state of the world into a repeatedly accessible signal.  The important
improvement is not necessarily that the world changes faster.  Instead,
\[
    I(F;T)
\]
increases, or informative \(F\) becomes available at shorter latency.

The same transformation occurs in engineering when an expensive full-system failure
is replaced by component tests, simulation, telemetry, or formal invariants.

\subsection{Verifier construction can reduce the feedback horizon in three ways}

The framework separates three distinct improvements.

\paragraph{Earlier observation.}
A new measurement can reveal a relevant effect before the final outcome occurs.
Continuous telemetry, biomarkers, in-situ spectroscopy, and precursor events are
examples.  If the earlier measurement is genuinely informative about the target, this
reduces effective \(\tauT\).

\paragraph{Cheaper repetition.}
Automation can make the same target-relevant experiment repeatable at much higher
throughput.  A-Lab integrates planning, synthesis, and characterization and executed
353 experiments in 17 days \parencite{szymanski2023alab}.  AlphaFlow similarly couples
experiment selection with in-situ measurement
\parencite{volk2023alphaflow}.  Neither system removes chemical or physical latency, but
both remove substantial human handoff latency and increase the number of
target-relevant observations per unit time.

\paragraph{Decomposition.}
A long-horizon target can sometimes be split into shorter independently checkable
conditions.  Software tests provide a familiar example.  A requirement such as
``this service remains correct in deployment'' is not directly verifiable in one short
experiment, but type checks, unit tests, property tests, integration tests, security
tests, and monitoring each constrain different failure modes.

In the ideal case,
\[
    T \approx g(T_1,\ldots,T_m),
\]
where each \(T_i\) has a much shorter verification horizon than \(T\).  The value of
the decomposition depends on whether satisfying the pieces is sufficient for the
whole.  Missing interactions can recreate the original long-horizon problem.

\subsection{Surrogates are compressed verifiers}

A common way to shorten the feedback horizon is to replace a slow target \(T\) with an
earlier observable \(P\):
\[
    T
    \longrightarrow
    P
    \longrightarrow
    \text{fast optimization}.
\]

This is verifier construction too, but it creates the central risk studied in this
paper.  A surrogate trades latency for approximation error.

Clinical medicine makes the trade explicit.  Accelerated approval can rely on
surrogate endpoints that are expected to predict clinical benefit, allowing decisions
before survival or quality-of-life evidence fully matures
\parencite{gyawali2019accelerated,naci2017accelerated}.  Later studies sometimes fail to
show corresponding patient-centered benefits
\parencite{liu2024clinicalbenefit}.  The surrogate has successfully shortened the
decision loop while imperfectly shortening the \emph{target} loop.

Reward models are the analogous construction in preference optimization.  Expensive
human judgments are compressed into a model that can be queried at machine speed.
This creates enormous evaluation throughput but also creates an exploitable residual:
continued optimization can increase proxy reward while independent evaluation
degrades \parencite{gao2022overoptimization,rafailov2024direct}.

Verifier construction therefore creates a trade:
\[
    \text{latency reduction}
    \quad\text{versus}\quad
    \text{loss of target information}.
\]

\subsection{A verifier is useful only relative to a claim}

There is no universal verifier for an artifact.

Lean verifies formal validity, not that the theorem is the one the human intended.
A unit test verifies tested behavior, not global program correctness.
A vulnerability reproduction harness verifies the presence of one exploit, not the
absence of others.
A biomarker measures a biological response, not necessarily patient benefit.

We can express this as a chain:
\[
\boxed{
\text{target specification}
\rightarrow
\text{observable}
\rightarrow
\text{verifier}
\rightarrow
\text{optimization}.
}
\]

Failure can occur at each arrow.

This also separates \emph{certification} from \emph{construction}.  Given a candidate
\(x\), an inexpensive verifier may answer whether \(x\) satisfies a criterion without
providing an efficient method for finding such an \(x\).  Formal theorem proving
makes the distinction especially clear: proof checking is cheap compared with proof
search.  Verifier construction improves the selection landscape; it does not by
itself solve the search problem.

\subsection{Optimization changes the requirements on the verifier}

A verifier that works for randomly sampled candidates need not remain reliable on
candidates selected specifically to maximize it.  Let \(P_0(x)\) be the distribution on
which an evaluator was validated and \(P_k(x)\) the distribution after \(k\) rounds of
optimization.  What matters is not only
\[
    \operatorname{Corr}_{P_0}(F,T),
\]
but whether the relationship survives the shift
\[
    P_0
    \rightarrow
    P_1
    \rightarrow
    \cdots
    \rightarrow
    P_K.
\]

Reward-model overoptimization provides direct evidence of this problem
\parencite{gao2022overoptimization}.  Continued policy optimization moves samples into
regions where the learned evaluator is less reliable.  More generally, optimization
searches for the residual error of the verifier.

A useful verifier for a powerful optimizer must therefore satisfy a stronger
criterion than ordinary predictive accuracy:
\[
    \text{high fidelity under the distribution induced by optimizing the verifier}.
\]

This requirement becomes especially difficult when the optimizer can influence the
measurement process itself.

\subsection{Independent correction is part of verifier design}

Verifier construction should therefore not aim only at making a single evaluator
faster.  It can also construct \emph{independent} correction channels.

Suppose a fast evaluator \(F_1\) guides most optimization, while evaluators
\(F_2,\ldots,F_m\) are queried less frequently.  The secondary evaluators are valuable
when their errors are not strongly correlated:
\[
    \Cov
    \bigl(
        \epsilon_i,\epsilon_j
    \bigr)
    \text{ is small},
    \qquad
    \epsilon_i = F_i-T.
\]

A second evaluator produced from the same data, architecture, and objective may add
little independent information even if its outputs look different.  Independence is
therefore a causal property of the evaluation process, not merely an ensemble size.

Human institutions often approximate this architecture through replication,
independent audits, adversarial review, appeals, and multiple measurement methods.
These mechanisms are imperfect, but their structural purpose is relevant here: they
prevent the process being optimized from relying indefinitely on one unchecked error
channel.

\subsection{Verifier construction is the main escape from a growing horizon gap}

If AI reduces \(\tauG\) and \(\tauP\) while \(\tauT\) remains fixed, then \(K\)
grows.  One response is to slow optimization.  A more powerful response, when
possible, is to redesign the evidence channel:
\[
\begin{array}{ccc}
\tauG \downarrow
&
\tauP \downarrow
&
\tauT \text{ fixed}
\\[3pt]
&\Downarrow&
\\[-2pt]
&K\uparrow&
\end{array}
\]
versus
\[
\begin{array}{ccc}
\tauG \downarrow
&
\tauP \downarrow
&
\tauT \downarrow
\\[3pt]
&\Downarrow&
\\[-2pt]
&K\text{ controlled}.&
\end{array}
\]

This gives a different interpretation of scientific and engineering progress.
Capability improvements search the solution space faster.  Measurement,
instrumentation, testing, and experimental design make the consequences of that
search legible faster.

The two are complements.  Fast search with a strong verifier produces the favorable
regime represented by theorem proving and well-tested software.  Fast search with a
weak but rapidly computed proxy produces the regime represented by reward-model
overoptimization.

For AI safety, the practical question is therefore not merely whether an evaluation
suite exists.  It is:

\begin{quote}
Does the evaluation architecture make target-relevant errors visible at roughly the
rate at which optimization can discover and amplify them?
\end{quote}

If not, increasing evaluator throughput alone may make the proxy loop faster without
making the correction loop faster.  The feedback horizon has not actually been
closed.
